\documentclass[svgnames]{article}    

\usepackage{graphicx}                 

% -- math                               
\usepackage{amssymb}
\usepackage{amsmath}
\usepackage{esint}
\usepackage{geometry}
\usepackage{dsfont}

% -- noindent
\setlength\parindent{0pt}

% -- images
%\graphicspath{{ }}                   

% -- tikz
\usepackage{pgfplots}
\pgfplotsset{compat=1.15}
\usepackage[most]{tcolorbox}

\newtcolorbox{greenbox}{
    enhanced,
    boxrule=0pt,frame hidden,
    colback=orange!10!white,
    sharp corners
}
\newtcolorbox{redbox}{
    enhanced,
    boxrule=0pt,frame hidden,
    colback=red!10!white,
    sharp corners
}
\newtcolorbox{bluebox}{
    enhanced,
    boxrule=0pt,frame hidden,
    colback=green!10!white,
    sharp corners
}
\newtcolorbox{purplebox}{
    enhanced,
    boxrule=0pt,frame hidden,
    colback=blue!10!white,
    sharp corners
}
\newtcolorbox{pinkbox}{
    enhanced,
    boxrule=0pt,frame hidden,
    colback=purple!10!white,
    sharp corners
}
\newtcolorbox{blackbox}{
    enhanced,
    colback=black!10!white,
    sharp corners
}


% -- figures
\usepackage{float}
\usepackage{caption}
\usepackage{lipsum}
\usepackage{comment}

% -- start 
\title{Deriving 4H-Silicon Carbide's Spin Hamiltonian}
\author{Deval Deliwala}
%\date{}                              

\begin{document}
\maketitle
%\tableofcontents  

The Spin Hamiltonian is a combination of the Zeeman effect, Hyperfine
Interactions, the Zero-Field Splitting Effect, and the Exchange Interaction. 

\begin{redbox}
  \[
      \hat{H}_0 = \hat{H}_Z + \hat{H}_{HF} + \hat{H}_{ZFS} + \hat{H}_{EX}
  \] 
\end{redbox}

Because Silicon Carbide (SiC) has no nuclei with nuclear spin quantum number $I
> \frac{1}{2}$, there is no Nuclear Quadrupole Interaction Hamiltonian. \\ 

The goal of this paper is to provide a thorough overview and derivation of each Hamiltonian
term. Finally, a complete Hamiltonian description will be provided along with
all unknown constants that need to be derived from experimental data. 

\tableofcontents

\newpage

\section{The Zeeman Effect}

A spinning charged particle is a magnetic dipole. Its magnetic dipole moment,
$\vec{\mu}$, is proportional to its spin angular momentum, $\vec{S}$: 
\[
\vec{\mu} = \gamma \vec{S}; 
\] 

the proportionality constant, $\gamma$, is the gyromagnetic ratio. From the
Dirac equation it can be shown  

\begin{align} \label{Dirac}
\vec{\mu} = \gamma \vec{S} = -g\frac{q}{2m_e}\vec{S} =
-\frac{g\mu_B}{\hbar}\vec{S}
\end{align} \vspace{3px}

where $g \approx 2.0023$ is the Lande $g$ factor and $\mu_B \approx 9.2 \cdot
10^{-24} J/T$ is the Bohr Magneton. \\

When a magnetic dipole is placed in a magnetic field $\vec{B}$, it
experiences a torque, $\vec{\mu}\times\vec{B}$, which tends to line it up
parallel to the field like a compass. The energy associated with this torque is 

\begin{align} \label{torque}
H = -\vec{\mu} \cdot \vec{B} = -\gamma \vec{B} \cdot S 
\end{align} 

Therfore, the Zeeman Hamiltonian, 

\[
\hat{H}_Z = -\mu_0 \cdot \vec{B}_0 = \frac{g\mu_B}{\hbar} \vec{S} \cdot
\vec{B}_0 = \frac{g\mu_B}{\hbar}B_0 S_z \] 

where $S_z = \frac{\hbar}{2}\hat{\sigma}_z$ and $\sigma_z$ is the Pauli-z spin matrix. Applying $\hat{H}_Z$ on the state $|s,
m_s\rangle$, we find 

\begin{align*}
    \hat{H}_z |s, m_s\rangle &= \frac{g\mu_B}{\hbar}B_0 \cdot S_z|s, m_s\rangle
                          \\ &= \frac{g\mu_B}{\hbar}B_0 \cdot m_s\hbar |s, m_s\rangle
                          \\ &= m_s g\mu_B B_0 |s, m_s\rangle.   
\end{align*} \vspace{3px} 

To generalize to an anisotropic $g$, we replace with a $3\times 3$ $g$ tensor, 

\[
\hat{H}_Z = \mu_B \vec{S} \cdot g \cdot \vec{B}_0
\] 

$g$ is symmetric and diagonalizable via the defect's \textit{principle axis}. 

\[
    g_{\text{diag}} = \begin{pmatrix}
        g_x & 0 & 0 \\ 
        0 & g_y & 0 \\ 
        0 & 0 & g_z 
    \end{pmatrix}  
\] 

However, for fitting purposes, setting $g\approx g_{\text{free } e^-}=2.0023$ is
acceptable. Hence this Hamiltonian is simple: 

\vspace{5px}
\begin{greenbox}
    \[
      \hat{H}_Z |s, m_s\rangle = m_s g\mu_B B_0 |s, m_s\rangle.
  \]
\end{greenbox}


\section{Hyperfine Interaction}

\subsection{Vector Potential $\vec{A}$}

From Maxwell's equations, we have 

\[
    \nabla \cdot \vec{B} = 0, \quad \nabla \times B = \mu_0 \vec{J}.
\] \vspace{3px}

Because $\nabla \cdot B = 0$, we set $\vec{B} = \nabla \times \vec{A}$. We
impose the Coulomb gauge $\nabla \cdot \vec{A} = 0$. Then 

\begin{align} \label{A}
\nabla^2 \vec{A} = -\mu_0 \vec{J}
\end{align}

\subsection{Current density $\vec{J}$}

We model the localized nuclear magnetic moment $\vec{\mu}_n$ at the origin by an
effective current density

\[
    \vec{J}_n (\vec{r}') = \nabla' \times [\vec{\mu}_n \delta^3(\vec{r}')]. 
\] \vspace{3px}

The general solution for A (Equation \ref{A}) is 

\[
\vec{A}(\vec{r}) = \frac{\mu_0}{4\pi} \int d^3 r'
\frac{\vec{J}_n(\vec{r}')}{|\vec{r}-\vec{r'}|} = \frac{\mu_0}{4\pi}\vec{\mu}_n
\times \nabla\left( \frac{1}{r} \right) = \frac{\mu_0}{4\pi}\frac{\vec{\mu}_n
\times \vec{r}}{r^3}.   
\] 

Therefore, 
\begin{align} \label{}
    \vec{B}(r) = \nabla \times \vec{A} = \frac{\mu_0}{4\pi}\nabla \times \left(
    \frac{\vec{\mu}_n \times \vec{r}}{r^3}\right).   
\end{align}

Using the identity 

\[
\nabla \times (\vec{a} \times \vec{r}f(r)) = \vec{a}\nabla \cdot (\vec{r}f) -
(\vec{a}\cdot \nabla)(\vec{r}f), 
\] \vspace{3px}

with $\vec{a} = \vec{\mu}_n$ and $f(r) = 1/r^3$, we get the canonical expression
for the magnetic field from a dipole from classical electrodynamics: 
\begin{align} \label{B}
      \vec{B}(\vec{r}) = \frac{\mu_0}{4\pi r^3} \left[ 3(\vec{\mu}_n \cdot
      \hat{r})\hat{r} - \vec{\mu}_n \right] + \frac{2\mu_0}{3}\vec{\mu}_n
      \delta^3(\vec{r}). 
\end{align}

\subsection{Hamiltonian Density}

A magnetic moment $\vec{\mu}_e$ at $\vec{r}$ has energy (Equation \ref{torque})

\[
U(\vec{r}) = -\vec{\mu}_e \cdot \vec{B}(\vec{r}). 
\] 

Therefore, from Equation \ref{B} we separate both terms, 
\begin{align*}
    U_\text{dipolar}(\vec{r}) &= -\vec{\mu}_e \cdot \left[\frac{\mu_0}{4\pi r^3}
    (3(\vec{\mu}_n \cdot \hat{r}) \hat{r} - \vec{\mu}_n)\right], \\
        U_\text{contact}(\vec{r}) &= -\vec{\mu}_e \cdot \left[
        \frac{2\mu_0}{3}\vec{\mu}_n \delta^3(\vec{r}) \right]. 
\end{align*}

We promote the magnetic moments to their quantum operators: 
\begin{align*}
    \vec{\mu}_e &= -g_e \mu_B \frac{\hat{S}}{\hbar}, \quad \mu_B =
    \frac{e\hbar}{2m_e}, \\
    \vec{\mu}_n &= g_n \mu_N \frac{\hat{I}}{\hbar}, \quad \mu_N =
    \frac{e\hbar}{2m_p}. 
\end{align*}

Therefore, the Hamiltonian of the electron, in the magnetic field due to the
nuclei's magnetic dipole moment at $\vec{r}$ (Equation \ref{B}), is $H = -\vec{\mu} \cdot \vec{B}$,
\begin{align} \label{HF_Density}
    \hat{H}(\hat{r}) = \frac{\mu_0}{4\pi} \frac{g_e\mu_B g_n \mu_N}{\hbar^2} \left(
    \frac{\vec{S} \cdot \vec{I} - 3(\vec{S} \cdot \hat{r})(\vec{I}\cdot
\hat{r})}{r^3} + \frac{8\pi}{3}\delta^3(\vec{r})\vec{S}\cdot \vec{I}\right).   
\end{align}


If the electron occupies spatial wave function $\psi(\vec{r})$, the full
hyperfine Hamiltonian is 

\[
    \hat{H}_{HF} = \int d^3r |\psi(\vec{r})|^2 \vec{H}(\vec{r}).
\] 

To solve this, we have 
\[ \hat{H}(\vec{r}) = C\left[ \frac{8\pi}{3}
        \delta^3(\vec{r})(\vec{S}\cdot\vec{I}) +
        \frac{1}{r^3}(\vec{S}\cdot\vec{I} - 3(\vec{S}\cdot
    \hat{r})(\vec{I}\cdot \hat{r}))\right], 
\]
where 

\[
C = \frac{\mu_0}{4\pi}\frac{g_e\mu_B g_n \mu_N}{\hbar^2}, \quad \hat{r}_i =
\frac{r_i}{r}. 
\]\vspace{3px} 

Writing the dipolar piece in index form, 

\begin{align} \label{index}
    \hat{H}(\vec{r}) = C \left[ \frac{8\pi}{3} \delta^3(\vec{r}) S_kI_k +
    \frac{1}{r^3} (\delta_{k\ell} - 3 \hat{r}_k\hat{r}_\ell) S_k I_\ell\right],
\end{align} \vspace{3px}

where we sum over repeated indices. Now to integrate, we have

\[
    \hat{H}_{HF} = \int d^3r |\psi(\vec{r})|^2 \vec{H}(\vec{r}) = \int d^3 r
    \rho(\vec{r}) \hat{H}(\vec{r}).
\] \vspace{3px}

where $\rho(r) = |\psi(\vec{r})|^2$ is the spatial probability density.
Splitting Equation \ref{index} into contact and dipolar integrals: 

\[
    \hat{H}_{HF} = C\left[ \frac{8\pi}{3}S_k I_k \int d^3r \rho(\vec{r})
        \delta^3(\vec{r}) \quad +\quad S_k I_\ell \int d^3r \rho(\vec{r})
    \frac{\delta_{k\ell} - 3 \hat{r}_k \hat{r}_\ell}{r^3}\right]. 
\] \vspace{3px}

The contact term (isotropic) reduces to 

\begin{align*}
    \frac{8\pi}{3}S_k I_k \int d^3 r \rho(\vec{r}) \delta^3(r) &=
\frac{8\pi}{3}S_kI_k \rho(0) \\ &= \frac{8\pi}{3}|\psi(0)|^2 S_k I_k \\ 
                                &= A_\text{iso}\delta_{k\ell} S_k I_\ell, \qquad A_\text{iso} =
C\frac{8\pi}{3}|\psi(0)|^2.    
\end{align*} 

To solve the dipolar term, we define 

\begin{align} \label{dipolar_}
    A_{k\ell}^\text{dip} = C\int d^3r \rho(\vec{r}) \frac{\delta_{k\ell} - 3
    \hat{r}_k \hat{r}_\ell}{r^3}. 
\end{align} \vspace{3px}

Then the term reduces to $S_kI_\ell A_{k\ell}^\text{dip}$. We define the
$A_{ij}$ tensor as 

\[
    A_{ij} = A_\text{iso} \delta_{ij} + A_{ij}^\text{dip}.
\] 

Therefore

\vspace{5px}
\begin{bluebox}
    \[ \hat{H}_{HF} = \hat{S} \cdot A \cdot \hat{I} = \hat{S}_i A_{ij} \hat{I}_j.
    \]
\end{bluebox}
\vspace{5px}

For multiple nuclei, like in SiC, we sum over every $j$ inequivalent nucleus. 

\[ \hat{H}_{HF} = \sum_j \hat{\vec{S}} \cdot A_j \cdot \hat{\vec{I}}_j.\] 

The dipolar integral in $A_{ij}^\text{dip}$ is often evaluated by choosing the
defect's principal axis frame to diagonalize. \\

From Equation \ref{spinxspiny}, the Cartesian spin operators are shown to be the
sum or difference of the raising or lowering operators. This relation also holds
true for Cartesian nuclear spin operators acting on nuclear spin. Therefore, with two nuclei, like
in SiC, the Hyperfine Hamiltonian acting on a state $|s, m\rangle
|m_{I_1}\rangle |m_{I_2}\rangle$ can be written as 

\section{Zero-Field Splitting}

The Zero-Field Splitting Hamiltonian is the dipole-dipole interaction between a
paramagnetic defect electron and the carrier electron. This is incredibly
similar to the Hyperfine interaction between a nearby nucleus and the carrier
electron. We can start from Equation \ref{B}

\[ 
      \vec{B}(\vec{r}) = \frac{\mu_0}{4\pi r^3} \left[ 3(\vec{\mu}_n \cdot
      \hat{r})\hat{r} - \vec{\mu}_n \right] + \left( \frac{2\mu_0}{3}\vec{\mu}_n
      \delta^3(\vec{r}) \right). 
  \]

However, here we remove the second term in the large parenthesis. This was the
contact ($\delta$-function) term only for a pointlike nucleus -- in Hyperfine,
the nucleus sits at a point, so when we take the curl of $\vec{A}$, we pick up
the singular term $\frac{2\mu_0}{3}\vec{\mu}_n\delta^3(\vec{r)}$. \\ 

However, for electron-electron interactions, two electrons never coincide
exactly, so there is no $\delta$-contact term. Only the $1/r^3$ dipolar piece
survives from taking the curl. Hence, we have 
\begin{align} 
      \vec{B_1}(\vec{r}) = \frac{\mu_0}{4\pi r^3} \left[ 3(\vec{\mu}_1 \cdot
      \hat{r})\hat{r} - \vec{\mu}_1 \right], \quad \hat{r} = \frac{\vec{r}}{r},  
\end{align}

where $\vec{\mu}_1$ is the magnetic moment of electron 1 and 
$\vec{B}_1(\vec{r})$ is the magnetic field of electron 1 at point
$\vec{r}$. A second dipole $\vec{\mu}_2$ located at $\vec{r}$ has energy in the
field of dipole 1: 

\[
U = -\vec{\mu}_2 \cdot \vec{B}_1(\vec{r}) = \frac{\mu_0}{4\pi r^3}\left[-\mu_1
\cdot \mu_2 +  3(\mu_1 \cdot \hat{r}) (\mu_2 \cdot \hat{r}) \right]. 
\] \vspace{3px}

Now we can promote $U \rightarrow \hat{H}_{ZFS}$ by using the electron magnetic
moment operator (Equation \ref{Dirac}): 

\[
    \hat{\vec{\mu}}_i = -g_e \mu_B \frac{\hat{\vec{S}}_i}{\hbar}, \quad \mu_B =
    \frac{e\hbar}{2m_e}, \; g_e \approx 2.0023. 
\] \vspace{3px}

Therefore, replacing $\vec{\mu}_{1,2}$ by $\hat{\vec{\mu}}_{1,2}$, the prefactor
becomes  
\[ -\frac{\mu_0}{4\pi} \frac{(g_e \mu_B)^2}{\hbar^2}. \] 

Therefore, we have 
\begin{align*}
    \hat{H}_{ZFS} &= -\frac{\mu_0}{4\pi}\frac{(g_e\mu_B)^2}{\hbar^2}\frac{1}{r^3}
    \left[ -\hat{\vec{S}}_1 \cdot \hat{\vec{S}}_2 + 3(\hat{\vec{S}}_1 \cdot
    \hat{r}) (\hat{\vec{S}}_2 \cdot \hat{r}) \right]. \\ 
                  &= \frac{\mu_0}{4\pi}\frac{(g_e \mu_B)^2}{\hbar^2}
                  \frac{1}{r^3}\left[\hat{\vec{S}}_1 \cdot \hat{\vec{S}}_2 -
                      3(\hat{\vec{S}}_1 \cdot \hat{r})(\hat{\vec{S}}_2 \cdot
                  \hat{r})\right]. 
\end{align*} \vspace{3px}

We can simplify this into an analogous form of the Hyperfine Hamiltonian. 
We have the equalities
\begin{align*}
    (\vec{S}_1 \cdot \vec{S}_2) &= \sum_{i\in \{x, y, z\}} S_{1i}S_{2i} \\ 
    (\vec{S}_1 \cdot \hat{r})(\vec{S}_2 \cdot \hat{r}) &= \left(\sum_i
    S_{1i}\hat{r_i}\right)\left(\sum_j S_{2j}\hat{r_j}\right)\\ &= \sum_{i,j}S_{1i}(\hat{r}_i
    \hat{r}_j)S_{2j}\\ &= \sum_{i,j} S_{1i}\frac{r_ir_j}{r^2}S_{2i}. 
\end{align*}

Hence, 
\begin{align*}
    \hat{H}_{ZFS} &= \frac{\mu_0}{4\pi}\frac{(g_e\mu_B)^2}{\hbar^2}
    \frac{1}{r^3} \left( \sum_{i}^{} S_{1i}S_{2i} -
    3\sum_{i,j}S_{1i}\frac{r_ir_j}{r^2}S_{2j} \right) \\
                  &= \frac{\mu_0}{4\pi}\frac{(g_e\mu_B)^2}{\hbar^2}\left( \sum_i
                  \frac{\delta_{i,j}}{r^3}S_{1i}S_{2j} - 3\sum_{i,j}
              \frac{r_ir_j}{r^5}S_{1i}S_{2j}\right) \\ 
                  &= \sum_{i,j} S_{1i} \left[ \frac{\mu_0}{4\pi}
                  \frac{(g_e\mu_B)^2}{\hbar^2} \left( \frac{\delta_{i,j}}{r^3} -
              \frac{3r_ir_j}{r^5}\right) \right] S_{2j}.  
\end{align*}

We define the tensor $D_{ij}$ 
\[
    D_{ij} (\vec{r}) = \frac{\mu_0}{4\pi} \frac{(g_e\mu_B)^2}{\hbar^2} \left(
    \frac{\delta_{ij}}{r^3} - \frac{3r_ir_j}{r^5} \right).   
\] \vspace{3px}

Then we get a similar equation to Hyperfine Hamiltonian: 
    \begin{align} 
        \hat{H}_{ZFS} = \sum_{i,j \in \{x, y, z\}}^{}  S_{1i} D_{ij}(\vec{r})
        S_{2i} = \hat{\vec{S}} \cdot D \cdot \hat{\vec{S}}. 
    \end{align} 

We can also diagonalize D via the principle axis. Additionally, defining
$K_{ij}(\vec{r})$, 

\[
    D_{ij}(\vec{r}) = \frac{\mu_0}{4\pi} \frac{(g\mu_B)^2}{\hbar^2} \left(
    \frac{\delta_{ij}}{r^3} - \frac{3r_i r_j}{r^5} \right) = \frac{\mu_0}{4\pi}
    \frac{(g_e\mu_B)^2}{\hbar^2} K_{ij}(\vec{r}),
\] \vspace{3px}

and taking the trace, 

\[
    \sum_i K_{ii} = \sum \frac{1}{r^3} - 3\sum \frac{r_i^2}{r^5} = \frac{3}{r^3} -
\frac{3}{r^2}(r_x^2 + r_y^2 + r_z^2) = 0. 
\] \vspace{3px}

Therefore, $D_{ij}$ is traceless. So if we diagonalize via the principle axis,
we have 

 \[
     D_{ij} = \begin{pmatrix}
    D_x & 0 & 0 \\ 
    0 & D_y & 0 \\ 
    0 & 0 & D_z 
\end{pmatrix}, \qquad D_x + D_y + D_z = 0. 
\] \vspace{3px}

And because the Silicon Vacancy $V_\text{Si}^- \in C_{3V}$ point group, we
have $D_x = D_y$. We can parameterize to fewer unknown variables.  

\[
\begin{cases}
    2D_x + D_z = 0 \\
               D_x = D_y 
\end{cases} \quad \Rightarrow \quad D_x = D_y = -\frac{1}{2}D_z. 
\] \vspace{3px}

Define $D, E$ such that 

 \[
D = \frac{3}{2}D_z \qquad E = \frac{1}{2}(D_x - D_y) = 0. 
\]

Then, we have 

\[
D_\text{diag} = \begin{pmatrix}
    D_x & 0 & 0 \\ 
    0 & D_y & 0 \\ 
    0 & 0 & D_z 
\end{pmatrix}  = \begin{pmatrix}
    -\frac{D}{3} + E & 0 & 0 \\ 
    0 & -\frac{D}{3} - E & 0 \\ 
    0 & 0 & \frac{2}{3}D 
\end{pmatrix}.
\] \vspace{3px}


However, more accurately, we integrate. $D$ is analogous to $A$ in Equation
\ref{dipolar_} and describes the dipole-dipole interaction between two electron
spins instead of electron-nuclear. We have 

\[
    D_{ij} = \frac{\mu_0}{4\pi} g_a g_b \mu_B^2 \left\langle \psi_0 \Bigg|
        \frac{3r_ir_j - \delta_{ij}r^2}{r^5} \Bigg| \psi_0\right\rangle,
\] \vspace{3px}

which is the exact same as Equation \ref{dipolar_}. Both Hamiltonians are
derived analogously. Thus,

\vspace{5px}
\begin{purplebox}
    \[ \hat{H}_{ZFS} = \hat{\vec{S}} \cdot D \cdot \hat{\vec{S}}.  
    \]
\end{purplebox}
\vspace{5px} 

We can expand, 
\begin{align}
    \hat{H}_{ZFS} &= \hat{\vec{S}}_a \cdot D \cdot \hat{\vec{S}}_b \\
                  &= \left( -\frac{D}{3}+ E \right)\hat{S}_x^2 - \left(
                  \frac{D}{3} + E \right) \hat{S}_y^2 + \frac{2}{3}D \hat{S}_z^2
                  \\ 
                  &= D\left(\hat{S}_z^2 - \frac{1}{3}\left(\hat{S}_x^2 +
                      \hat{S}_y^2 + \hat{S}_z^2\right) \right) + E(\hat{S}_x -
                      \hat{S}_y^2), 
\end{align}

With the $\hat{S}_x, \hat{S}_y$ Cartesian spin operators being related to the
raising and lowering operators via 

\begin{align} \label{raisinglowering}
    \hat{S}_+ |s, m\rangle &= \frac{1}{2}(\hat{S}_x + i \hat{S}_y )|s, m\rangle \\ 
                           &= \sqrt{(s(s+1) - m(m+1))}|s, m+1\rangle. \\ 
    \hat{S}_- |s, m\rangle &= \frac{1}{2}(\hat{S}_x - i \hat{S}_y )|s, m\rangle \\
                           &= \sqrt{(s(s+1) - m(m-1))}|s, m-1\rangle. 
\end{align}

By rearranging Equation \ref{raisinglowering} in terms of $\hat{S}_x,
\hat{S}_y$, we get 

\begin{align} \label{spinxspiny}
    \hat{S}_x &= \frac{1}{2}(\hat{S}_+ + \hat{S}_-) \\ 
    \hat{S}_y &= \frac{1}{2i}(\hat{S}_+ - \hat{S}_-) \\ 
    \hat{S}_x^2 + \hat{S}_y^2 = \hat{S}_x^2 - \hat{S}_y^2 &=
    \frac{1}{2}(\hat{S}_+^2 + \hat{S}_-^2).  
\end{align}

Therefore, a general expression for the Zero-Field Splitting Hamiltonian acting
on $|s, m\rangle$ can be written as 

\begin{purplebox}
  \begin{align*}
      \hat{H}_{ZFS} |s, m\rangle &= Dm^2 |s, m\rangle - \frac{D}{3}(s(s+1))|s,
      m\rangle \\ &+ \frac{E}{2}(s(s+1) - m(m+1))|s, m+2\rangle \\ 
                  &+ \frac{E}{2}(s(s+1) - m(m-1))|s, m-2\rangle. 
  \end{align*}
  \vspace{2px}
\end{purplebox}

\section{Exchange Interaction}

Consider two electrons in orbitals $\phi_a(\vec{r}), \phi_b(\vec{r})$. The
Hamiltonian is 

\[
\hat{H} = \sum_{i=1}^{2} \left[ -\frac{\hbar^2}{2m} \nabla_i^2 +
V(\vec{r}_i)\right] + \frac{e^2}{4\pi\varepsilon_0 |\vec{r}_1 - \vec{r}_2 |}. 
\] \vspace{3px}

Here, $V(\vec{r})$ is the single-particle potential. Because electrons are
fermions, their total state must be antisymmetric under spatial and spin
exchange. We factorize into spatial $\Psi(\vec{r}_1, \vec{r}_2)$ and spin
$\chi(s_1, s_2)$ parts: 

\[
\Psi_\text{tot} (1, 2) = \Psi(\vec{r}_1, \vec{r}_2) \chi(s_1, s_2),  
\] 

with 

\[ \Psi_\text{tot}(2, 1) = -\Psi_\text{tot}(1, 2).  \] 

If $\Psi$ is symmetric, then  $\chi$ must be antisymmetric (singlet). If $\Psi$
is antisymmetric, $\chi$ must be symmetric (triplet). \\ 

We build orthonormal spatial orbitals  $\phi_a, \phi_b$. Two natural
two-electron spatial states are: 
\begin{align*}
    \Psi_S(\vec{r}_1, \vec{r}_2) &= \frac{1}{\sqrt{2}}\left[\phi_a(1)\phi_b(2) +
    \phi_b(1) \phi_a(2) \right] \\ 
        \Psi_A(\vec{r}_1, \vec{r}_2) &= \frac{1}{\sqrt{2}} \left[ \phi_a(1)
        \phi_b(2) - \phi_b(1) \phi_a(2) \right]. 
\end{align*} 

The corresponding spin states are 
\begin{align*}
    \chi_S &= \frac{1}{\sqrt{2}} (|\uparrow \downarrow\rangle - |\downarrow
    \uparrow \rangle) \quad \text{(singlet)} \\  
    \chi_A &= \frac{1}{\sqrt{2}} (|\uparrow \downarrow\rangle + |\downarrow
    \uparrow \rangle) \quad \text{(triplet)}. 
\end{align*}

We diagonalize $\hat{H}$ into the two-dimensional subspace spanned by
$\phi_a\phi_b$. Because  $\hat{H}_0$ acts on each electron separately and both
spatial states $\Psi_S, \Psi_A$ have the same one-particle energy, 

\[
E_0 = \langle \phi_a | \hat{h} | \phi_a\rangle + \langle \phi_b | \hat{h} |
\phi_b\rangle, \qquad \hat{h} = -\frac{\hbar^2}{2m}\nabla^2 + V(\vec{r});
\] \vspace{3px}

their splitting arises entirely from the coulomb term: 

\[
    \hat{H}_C = \frac{e^2}{4\pi\varepsilon_0 r_{12}}
\] \vspace{3px}

Defining the direct and exchange integrals: 
\begin{align*}
    K &\equiv \int \int d^3r_1 d^3r_2 |\phi_a(1)|^2 \frac{e^2}{4\pi\varepsilon_0
    r_{12}} |\phi_b(2)|^2, \\ 
        J &\equiv \int\int d^3r_1 d^3r_2 \phi_a^*(1) \phi_b(1)
        \frac{e^2}{4\pi\varepsilon_0r_{12}} \phi_b^*(2) \phi_a(2). 
\end{align*}

We find 
\begin{align*}
    \langle \Psi_S | \hat{H} | \Psi_S \rangle &= E_0 + K + J, \\
    \langle \Psi_A | \hat{H} | \Psi_A \rangle &= E_0 + K - J. 
\end{align*}

Thus the singlet sits at energy $E_S = E_0 + K + J$, while each triplet has $E_T
= E_0 + K - J$. Their energy difference/splitting is

\[
\Delta E \equiv E_S - E_T = 2J. 
\] \vspace{3px}


The energy splitting due to $\hat{H}_{ZFS}$ is analogously also $2D$.
Regardless, we now wish to replace this two level spatial+spin system by a
purely spin Hamiltonian that reproduces the same energy shift between singlet
and triplet.  \\ 

Recall that the operator exchanging two spins is 

\[
    \hat{P}_{12} \chi(1, 2) = \chi(2, 1), 
\]

It's eigenvalues are +1 on the symmetric (triplet) subspace and -1 on the
antisymmetric (singlet) subspace. Define the total spin operator 

\[
\hat{\vec{S}}_\text{tot} = \hat{\vec{S}}_1 + \hat{\vec{S}}_2, 
\] 

and its square 

\[
\hat{S}^2_\text{tot} = (\hat{\vec{S}}_1^2 + \hat{\vec{S}}_2)^2 + \hat{S}_1^2 +
 \hat{S}_2^2 + 2 \hat{\vec{S}}_1 \cdot \hat{\vec{S}}_2.  
\] \vspace{3px}

For spin-$\frac{1}{2}$, $\hat{S}_1^2 = \hat{S}_2^2 = \hat{S}_x^2 + \hat{S}_y^2 +
\hat{S}_z^2 = \frac{3\hbar^2}{4}\mathds{1}$. Hence, 

\begin{align} \label{Stot}
    \hat{S}_\text{tot}^2 = \frac{3\hbar^2}{2}\mathds{1} + 2 \hat{\vec{S}}_1
    \cdot \hat{\vec{S}}_2. 
\end{align} \vspace{3px}

Now, for both triplet ($S_\text{tot}=1$) and singlet ($S_\text{tot}=0$) states, 
\begin{align*}
    \hat{S}_\text{tot}^2 |10\rangle  &= S(S+1)\hbar^2|10\rangle =
    2\hbar^2|10\rangle \quad &&\text{We want $\hat{P}_{12} |10\rangle =
    +|10\rangle$}\\
    \hat{S}_\text{tot}^2 |00\rangle &= S(S+1)\hbar^2|00\rangle = 0 \quad
                                    &&\text{We want $\hat{P}_{12}|00\rangle = -|00\rangle$} 
\end{align*}

We define an ansatz $\hat{P}_{12} = a \hat{S}_\text{tot}^2 + b\mathds{1}$. Hence 
\begin{align*}
    \hat{P}_{12} |10\rangle &= a(2) + b |10\rangle = +1 |10\rangle \\ 
    \hat{P}_{12} |00\rangle &= a(0) + b |00\rangle = -1 |00\rangle. 
\end{align*} 

Solving, 

\[
b = -1, \quad 2a-1=+1 \; \Rightarrow \; 2a=2 \; \Rightarrow \; a = 1. 
\] 

Therefore, 

\[
    \hat{P}_{12} = 1 \cdot \hat{S}_\text{tot}^2 - 1 \cdot \mathds{1} =
    \hat{S}_\text{tot}^2 - \mathds{1}.   
\] 


Substituting back in $\hat{S}_\text{tot}^2$ from Equation \ref{Stot}, 

\[
    \hat{P}_{12} = \left( \frac{3}{2}\mathds{1} + 2 \hat{\vec{S}}_1 +
    \hat{\vec{S}}_2 \right) - \mathds{1} = \frac{1}{2}\mathds{1} + 2
    \hat{\vec{S}}_1 \cdot \hat{\vec{S}}_2.   
\] \vspace{3px}


We get the final expression for the exchange operator: 


\[
    \hat{P}_{12} = \frac{1}{2}\left(1 + 4 \hat{\vec{S}}_1 \cdot \hat{\vec{S}}_2
    \right).
\] \vspace{3px}

It's eigenvalues are -1 on the singlet and +1 on the triplet manifold. Therefore
an operator proportional to $\hat{P}_{12}$ will split singlet/triplet. Define
the proportionality coefficient $J_{ex}$ such that 

 \[
     \hat{H}_{EX} = J_{ex} \hat{P}_{12} = \frac{J_{ex}}{2}\left(1 + 4 \hat{\vec{S_1}} \cdot
     \hat{\vec{S}}_2\right). 
\] \vspace{3px}

Therefore, acting on 
\begin{align*}
    \chi_S: \; \hat{H}_{EX} \chi_S &= J_{ex} (-1) \chi_S \\ 
    \chi_T: \; \hat{H}_{EX} \chi_T &= J_{ex} (+1) \chi_T. 
\end{align*}

To reproduce $\Delta E = E_S - E_T = 2J$, we set  $J_{ex} = -J$. Then, 

 \[
     E_S = -J_{ex} = +J, \qquad E_T = +J_{ex} = -J, 
\] \vspace{3px}

as desired. Setting $J$ to be the exchange couling matrix proportional to the
degree of overlap between wavefunctions, the exchange Hamiltoniain can be
expressed as 

 \[
     \hat{H}_{EX} = \hat{\vec{S}}_a \cdot J \cdot \hat{\vec{S}}_b. 
\] \vspace{3px}

If we assume the exchange coupling is isotropic, we can simplify to 

\vspace{5px}
\begin{pinkbox}
     \[   \hat{H}_{EX} = -J \hat{\vec{S}}_a \cdot \hat{\vec{S}}_b. 
     \]
\end{pinkbox}
\vspace{5px}

\section{Putting it all together}

For a two-electron (carrier + defect) and two-nuclei (Silicon and Carbon) spin
system, there are 16 orthonormal basis states that make up the 16 $\times$ 16
Spin Hamiltonian. Written in the coupled-electron  $|s, m\rangle$ basis tensored
with each nucleus  $m_I$ eigenstate: 

\subsubsection*{Triplet  $|1, +1\rangle$ manifold}
\begin{align*}
  |1, 1\rangle \otimes |+\frac{1}{2}\rangle_1 \otimes | +\frac{1}{2} \rangle_2 \\ 
  |1, 1\rangle \otimes |+\frac{1}{2}\rangle_1 \otimes | -\frac{1}{2} \rangle_2 \\ 
  |1, 1\rangle \otimes |-\frac{1}{2}\rangle_1 \otimes | +\frac{1}{2} \rangle_2 \\ 
  |1, 1\rangle \otimes |-\frac{1}{2}\rangle_1 \otimes | -\frac{1}{2} \rangle_2 \\ 
\end{align*}
\subsubsection*{Triplet  $|1, \;\;\,0\rangle$ manifold}
\begin{align*}
  |1, 0\rangle \otimes |+\frac{1}{2}\rangle_1 \otimes | +\frac{1}{2} \rangle_2 \\ 
  |1, 0\rangle \otimes |+\frac{1}{2}\rangle_1 \otimes | -\frac{1}{2} \rangle_2 \\ 
  |1, 0\rangle \otimes |-\frac{1}{2}\rangle_1 \otimes | +\frac{1}{2} \rangle_2 \\ 
  |1, 0\rangle \otimes |-\frac{1}{2}\rangle_1 \otimes | -\frac{1}{2} \rangle_2 \\ 
\end{align*}
\subsubsection*{Singlet  $|0, \;\;\,0\rangle$ manifold}
\begin{align*}
  |0, 0\rangle \otimes |+\frac{1}{2}\rangle_1 \otimes | +\frac{1}{2} \rangle_2 \\ 
  |0, 0\rangle \otimes |+\frac{1}{2}\rangle_1 \otimes | -\frac{1}{2} \rangle_2 \\ 
  |0, 0\rangle \otimes |-\frac{1}{2}\rangle_1 \otimes | +\frac{1}{2} \rangle_2 \\ 
  |0, 0\rangle \otimes |-\frac{1}{2}\rangle_1 \otimes | -\frac{1}{2} \rangle_2 \\ 
\end{align*}
\subsubsection*{Triplet  $|1, -1\rangle$ manifold}
\begin{align*}
  |1, -1\rangle \otimes |+\frac{1}{2}\rangle_1 \otimes | +\frac{1}{2} \rangle_2 \\ 
  |1, -1\rangle \otimes |+\frac{1}{2}\rangle_1 \otimes | -\frac{1}{2} \rangle_2 \\ 
  |1, -1\rangle \otimes |-\frac{1}{2}\rangle_1 \otimes | +\frac{1}{2} \rangle_2 \\ 
  |1, -1\rangle \otimes |-\frac{1}{2}\rangle_1 \otimes | -\frac{1}{2} \rangle_2 \\ 
\end{align*}

Our goal is to now programmatically calculate the Spin Hamiltonian's effect on
each coupled-basis state above. To do this, we calculate the Zeeman, Hyperfine,
ZFS, and Exchange Hamiltonian separately in the coupled basis, and then sum them
together. 
























\end{document}
